\documentclass[12px,a4paper]{article}
\usepackage[utf8]{inputenc}
\usepackage[portuguese]{babel}
\usepackage[T1]{fontenc}
\usepackage{graphicx}
\usepackage{xcolor}
\usepackage{titlesec}
\usepackage{geometry}
\geometry{a4paper, margin=1in}
\usepackage{hyperref}
\usepackage{tcolorbox}
\usepackage{enumitem}

% Definição de Cores Institucionais
\definecolor{CIECCRed}{HTML}{800000}
\definecolor{CIECCGold}{HTML}{C5A059}
\definecolor{CIECCNavy}{HTML}{111827}
\definecolor{CIECCCream}{HTML}{FDF5E6}

% Configuração de Títulos
\titleformat{\section}{\color{CIECCRed}\normalfont\Large\bfseries\scshape}{\thesection}{1em}{}
\titleformat{\subsection}{\color{CIECCNavy}\normalfont\large\bfseries}{\thesubsection}{1em}{}

\title{
    \vspace{2cm}
    \includegraphics[width=0.3\textwidth]{src/assets/logo-placeholder.svg} \\ % Placeholder para logo
    \vspace{1cm}
    \textcolor{CIECCRed}{\Huge \textbf{Proposta Técnica e Comercial}} \\
    \vspace{0.5cm}
    \textcolor{CIECCNavy}{\Large App Oficial do II CIECC 2026} \\
    \vspace{1cm}
    \large Hub Digital de Experiência e Engajamento
}
\author{Antigravity AI Coding Assistant}
\date{\today}

\begin{document}

\maketitle
\newpage

\section{Introdução e Objetivo}
O aplicativo oficial do \textbf{II Congresso Internacional de Educação Cristã Clássica (CIECC)} foi concebido para ser o epicentro da comunicação e relacionamento durante o evento. Mais do que um simples repositório de informações, ele atua como um \textit{Digital Concierge}, guiando o congressista em sua jornada acadêmica e espiritual.

\subsection{Proposta de Valor}
\begin{itemize}
    \item \textbf{Centralização:} Todas as informações (agenda, tickets, avisos) em um único lugar.
    * \textbf{Engajamento:} Ferramentas de networking para conectar educadores, gestores e famílias.
    * \textbf{Exclusividade:} Acesso a conteúdos de mídia e transmissões oficiais.
    * \textbf{Eficiência:} Redução de filas e otimização do fluxo de credenciamento via QR Code.
\end{itemize}

\section{Conceito Visual (Branding)}
A identidade visual do aplicativo é baseada no conceito \textbf{Classical Academic Premium}. Inspirada na tradição das artes liberais, utilizamos uma paleta de cores institucional que transmite autoridade e seriedade:
\begin{description}
    \item[Bordô (Crimson):] Simboliza a paixão pelo conhecimento e a herança acadêmica.
    \item[Dourado (Ochre):] Representa o valor eterno da sabedoria clássica.
    \item[Creme (Parchment):] Remete à tradição dos grandes textos e manuscritos.
\end{description}

\section{Arquitetura da Solução}
O aplicativo é dividido em cinco pilares fundamentais de navegação, desenhados com foco em UX \textit{mobile-first}:

\subsection{Dashboard de Início (Home)}
O painel de controle do congressista. Reúne banners de boas-vindas, contagem regressiva, comunicados de alta prioridade e atalhos rápidos para as funções mais utilizadas.

\subsection{Gestão de Programação (Agenda)}
Permite ao usuário visualizar a grade completa de atividades, filtrar por temas ou palestrantes e favoritar sessões de interesse para montar sua agenda personalizada.

\subsection{Networking e Relacionamento}
Área dedicada à conexão entre participantes. Facilita o encontro de perfis estratégicos (Gestores, Famílias Educadoras) e promove o relacionamento comercial com parceiros e expositores.

\section{Impacto e Resultados Esperados}
O app visa transformar a experiência do congresso em algo dinâmico e interativo, proporcionando:
\begin{description}
    \item[Para o Organizador:] Melhor controle de tráfego, comunicação instantânea e dados de interesse do público.
    \item[Para o Participante:] Sensação de valor institucional e facilidade na navegação pelo evento.
    \item[Para o Patrocinador:] Espaços nobres de visibilidade e conexão direta com o público-alvo.
\end{description}

\vspace{2cm}
\begin{tcolorbox}[colback=CIECCRed!5,colframe=CIECCRed,title=Nota de Desenvolvimento]
    Este documento é atualizado dinamicamente conforme a evolução do protótipo digital.
\end{tcolorbox}

\end{document}
